\section{Resumen ejecutivo}

GymHub es una plataforma web orientada a la gestión administrativa y al
seguimiento deportivo de gimnasios pequeños y medianos. El prototipo centraliza
miembros, membresías, pagos, asistencia, planes de entrenamiento, sesiones,
progreso, alertas y guías generales de nutrición. También incorpora paneles por
rol y un asistente conversacional contextual. La propuesta se plantea como un
servicio por suscripción y no como una venta única de software.

El proyecto parte de una oportunidad general documentada: las pequeñas y
medianas empresas suelen adoptar herramientas digitales a menor velocidad que
las organizaciones grandes, debido a restricciones de recursos, capacidades y
financiamiento \parencite{oecd2021}. Esta referencia respalda la pertinencia de
una solución sencilla y gradual, pero no demuestra por sí sola que todos los
gimnasios de Pérez Zeledón presenten el mismo problema. Por ese motivo, el
estudio local se mantiene pendiente hasta aplicar la encuesta incluida en los
anexos.

GymHub se encuentra en estado de MVP funcional y demostrable. El sistema tiene
una interfaz React, una API Django, persistencia PostgreSQL en Supabase y un
despliegue de demostración en Vercel. No obstante, aún requiere validación de
mercado, fortalecimiento de infraestructura, formalización de soporte y pruebas
con usuarios antes de presentarse como un servicio comercial maduro. El modelo
de negocio se formula con esa diferencia explícita entre prototipo técnico y
empresa operativa.

\section{Justificación}

\subsection{Problema y necesidad identificada}

La administración de un gimnasio exige mantener información actualizada sobre
personas usuarias, vigencia de membresías, pagos, ingreso a las instalaciones,
planes asignados y seguimiento del entrenamiento. Cuando esas tareas se
resuelven con registros separados, se incrementa la posibilidad de duplicar
datos, perder trazabilidad o dedicar tiempo a buscar información. La necesidad
que atiende GymHub consiste en disponer de una fuente centralizada que permita
consultar el estado administrativo y deportivo de cada miembro según su rol.

El documento inicial del proyecto identificó cuadernos, hojas de cálculo y
mensajería como métodos que se desea sustituir. Esta observación constituye la
hipótesis de problema del equipo, no un resultado estadístico. Su frecuencia y
gravedad en el mercado local deberán medirse con responsables de gimnasios y
miembros adultos. Hasta contar con esas respuestas, cualquier generalización
sobre Pérez Zeledón se marcará como pendiente de validar.

\subsection{Oportunidad detectada}

La digitalización puede reducir costos de transacción, integrar procesos y
facilitar el acceso de las pymes a información para tomar decisiones. Al mismo
tiempo, las empresas pequeñas enfrentan barreras de habilidades, tiempo y
financiamiento, por lo que una herramienta compleja puede resultar tan poco útil
como la ausencia de una solución \parencite{oecd2021}. GymHub busca ocupar ese
espacio mediante una interfaz responsive, incorporación guiada y crecimiento
por módulos.

La oportunidad comercial se concentra primero en dueños y administradores,
porque son quienes deciden la contratación. Entrenadores y miembros son usuarios
del sistema y determinan que la plataforma se use de manera constante. Esta
separación evita confundir al comprador con todos los beneficiarios.

\subsection{Contexto social, económico, tecnológico y ambiental}

En el ámbito social, una gestión ordenada puede ofrecer mayor claridad sobre
rutinas, asistencia y estado de la membresía. GymHub no garantiza resultados
físicos ni reemplaza el acompañamiento profesional; organiza información para
que las personas autorizadas puedan dar seguimiento.

En el ámbito económico, el valor esperado proviene del ahorro de tiempo, la
reducción de errores de registro y la visibilidad sobre pagos pendientes. Estos
beneficios deberán medirse durante un piloto, ya que todavía no se dispone de
datos para monetizar el ahorro producido.

En el ámbito tecnológico, el proyecto combina una aplicación web, una API, una
base de datos relacional, autenticación por roles, tareas programadas y servicios
de nube. La arquitectura permite demostrar el concepto sin comprar servidores,
pero los planes gratuitos actuales no deben confundirse con infraestructura
comercial con respaldo y soporte.

En el ámbito ambiental, el uso de registros digitales puede disminuir formularios
impresos. No se afirma un ahorro ambiental cuantificado porque el proyecto no ha
medido consumo de papel ni la huella de los servicios de nube. El beneficio se
presenta como potencial y sujeto a comprobación.

\subsection{Pertinencia del modelo de suscripción}

El servicio requiere alojamiento, mantenimiento, correcciones, copias de
seguridad y soporte continuo. Por ello, la suscripción mensual guarda mayor
relación con la estructura de costos que una venta única. Los planes escalonados
permiten iniciar con operación administrativa y habilitar funciones de mayor
valor conforme exista demanda. El precio definitivo dependerá de la disposición
de pago observada en la encuesta y de los costos reales del piloto.

\section{Objetivos del modelo de negocio}

\subsection{Objetivo general}

Formular el modelo de negocio de GymHub, una plataforma web para centralizar la
gestión administrativa y el seguimiento deportivo de gimnasios pequeños y
medianos de Pérez Zeledón, mediante el análisis de su propuesta de valor,
viabilidad técnica, costos y oportunidades comerciales durante 2026.

\subsection{Objetivos específicos}

\begin{enumerate}
  \item Caracterizar las necesidades de responsables y miembros de gimnasios
  mediante fuentes documentales y un instrumento de encuesta pendiente de
  aplicación.
  \item Documentar las funciones, arquitectura, seguridad y estado verificable
  del MVP de GymHub.
  \item Definir los nueve componentes del modelo Canvas y una estrategia de
  comercialización por suscripción.
  \item Estimar costos monetarios y aportes en especie, junto con escenarios
  provisionales de ingresos y punto de equilibrio.
  \item Establecer una ruta de validación, privacidad, soporte y mejora necesaria
  para una futura operación comercial.
\end{enumerate}
